% !TeX program = xelatex
\documentclass{homework}

\newcommand{\hwname}{XXX}
\newcommand{\hwid}{20XXXXXXXXXX}
\newcommand{\hwemail}{}
\newcommand{\hwclass}{示例课程}
\newcommand{\hwtype}{作业}
\newcommand{\hwnum}{1}

\begin{document}
\maketitle

\question\label{q:use}
本模板在 \href{https://github.com/jez/latex-homework-class}{\texttt{homework}} 类的题目接口上改为中文，版式取自国科大大作业模板：A4、小四、页边距 $2.5\,\mathrm{cm}$。页眉为课程与作业序号，默认每页带校徽水印。

个人信息写在导言区的 \verb|\hwname|、\verb|\hwid|、\verb|\hwclass| 等命令中。行内公式 $v-\varepsilon+\phi=2$，编号公式见\autoref{eq:euler}。

\begin{equation}
  v-\varepsilon+\phi=2
  \label{eq:euler}
\end{equation}

\question*{归纳：等差数列求和}\label{q:ind}
证明对任意正整数 $n$，
\[
  \sum_{i=1}^{n} i = \frac{n(n+1)}{2}.
\]

\begin{induction}
  \basecase
  当 $n=1$ 时，左边为 $1$，右边为 $1$，等式成立。
  \indhyp
  假设 $n=k$ 时成立，即 $\sum_{i=1}^{k} i = k(k+1)/2$。
  \indstep
  当 $n=k+1$ 时，
  \[
    \sum_{i=1}^{k+1} i
    = \frac{k(k+1)}{2}+(k+1)
    = \frac{(k+1)(k+2)}{2}.
  \]
\end{induction}

因此命题对所有正整数成立。本题编号为第\ref{q:ind}题。

\question
\verb|arabicparts| 的小问带当前题号。
\begin{arabicparts}
  \questionpart
  时间复杂度记为 $T(n)$。
  \questionpart
  证明 $T(n)=O(n\log n)$。
\end{arabicparts}

\question
\verb|alphaparts| 使用字母编号。同一题里分两段书写时，编号连续。
\begin{alphaparts}
  \questionpart
  打开终端。
  \questionpart
  执行 \texttt{xelatex homework.tex}。
\end{alphaparts}

中间可以插入说明，再继续小问：

\begin{alphaparts}
  \questionpart
  检查页眉、水印和题号。
\end{alphaparts}

嵌套列举仍然可用。
\begin{alphaparts}
  \questionpart
  编译顺序
  \begin{enumerate}
    \item \texttt{xelatex}
    \item 若加入参考文献，再跑 \texttt{bibtex}，然后 \texttt{xelatex} 两遍
  \end{enumerate}
  \questionpart
  或直接 \texttt{make}。
\end{alphaparts}

\renewcommand{\questiontype}{练习}
\question
\verb|\renewcommand{\questiontype}{练习}| 只改变编号题的前缀，这里是「练习」而不是「题目」。

\renewcommand{\questiontype}{题目}

\setcounter{questionCounter}{9}
\question
\verb|\setcounter{questionCounter}{9}| 使下一题从 10 开始。

\section
无星号的 \verb|\section| 与 \verb|\question| 相同，适合从 Markdown 转来的稿件。

\section*{使用 section 星号}
\verb|\section*{标题}| 与 \verb|\question*{标题}| 相同。

\question
代码、表格、图片、定理和作答框。

\begin{lstlisting}[language=Python]
def prefix_sum(values):
    total = 0
    output = []
    for value in values:
        total += value
        output.append(total)
    return output
\end{lstlisting}

\begin{table}[!htbp]
  \centering
  \begin{tabular}{@{}ll@{}}
    \toprule
    文件 & 作用 \\
    \midrule
    \texttt{homework.cls} & 文档类 \\
    \texttt{template.tex} & 空白起点 \\
    \texttt{homework.tex} & 本示例 \\
    \texttt{figures/} & 校徽 \\
    \bottomrule
  \end{tabular}
  \caption{模板文件}
  \label{tab:files}
\end{table}

\autoref{tab:files}给出文件组成。校徽见\autoref{fig:logo}。

\begin{figure}[!htbp]
  \centering
  \includegraphics[width=0.46\textwidth]{ucas_logo.pdf}
  \caption{中国科学院大学}
  \label{fig:logo}
\end{figure}

\begin{definition}
  设 $f,g:\mathbb{N}\to\mathbb{R}_{\geq 0}$。若存在常数 $c>0$ 和 $n_0$，使得对所有 $n\geq n_0$ 有 $f(n)\leq c\,g(n)$，则称 $f(n)=O(g(n))$。
\end{definition}

\begin{theorem}
  若 $f(n)=O(n)$ 且 $g(n)=O(n)$，则 $f(n)+g(n)=O(n)$。
\end{theorem}

\begin{proof}
  取两种上界的常数之和，阈值取两者的较大者。
\end{proof}

\tbox{默认开启水印。\texttt{\textbackslash documentclass[nowatermark]\{homework\}} 关闭水印，并改为偶数页页眉印姓名、奇数页页眉印校徽。匿名提交使用 \texttt{anonymous}，每题新页使用 \texttt{newpage}。}

\begin{samepage}
手写作答可留空框：
\answerbox{2.6cm}
\end{samepage}

\end{document}
